% Requires \usepackage{booktabs} and the bibliography entries
% liu2026loopahead and chen2026dflash in suffix-loop-design.bib.
% Scope: the implemented Ouro-2.6B and Huginn-3.5B training protocols.
% Huginn 332.6M projected DFlash training was ongoing when this section was prepared;
% no final training cost or completed-run claim is included.

\subsection{Training the LoopAhead and DFlash Baselines}
\label{app:drafter-training}

We follow the architecture and training recipes reported for
LoopAhead~\cite{liu2026loopahead} and DFlash~\cite{chen2026dflash},
with the data and target-model adaptations described below.
We train a separate drafter for each target model, keeping the target
parameters, token embeddings, and output projection frozen.
Our targets are Ouro-2.6B with four recurrent passes and Huginn-3.5B
with 32 passes. Huginn is an adaptation beyond the target models
evaluated in the original papers.

\paragraph{Training data.}
Table~\ref{tab:drafter-data} summarizes the prompt pools.
For LoopAhead, we retain the original split sizes and domain quotas,
but replace its BBH component with Nemotron V1 STEM questions.
The original BBH split identifiers are unavailable, so this substitution
allows us to exclude the entire BBH evaluation source from training.
For DFlash, we use a predefined budget of 50,000 prompts, compared with
approximately 800,000 examples in the original paper, allocating 40\%
to mathematics, 30\% to code, and 30\% to dialogue.
We use the publicly accessible Nemotron V1 because V2 was access-restricted
during data preparation. Missing Nemotron code prompts are recovered
from verified TACO training-set records, and UltraChat replaces the
unavailable dialogue prompts. These are adaptations of the published
data recipes, rather than identical training corpora.

\begin{table*}[t]
  \centering
  \small
  \caption{Prompt counts for training and held-out splits.
  Counts precede target-specific response filtering.
  UltraChat uses the first user request from its \texttt{train\_sft} split.
  The confirmation split is used only for LoopAhead.}
  \label{tab:drafter-data}
  \begin{tabular}{llrrr}
    \toprule
    Method & Source & Train & Development & Confirmation \\
    \midrule
    LoopAhead & GSM8K training split & 2,000 & 128 & 128 \\
              & CodeAlpaca, Python subset & 2,000 & 128 & 128 \\
              & MATH training split & 2,000 & 128 & 128 \\
              & Nemotron V1 STEM & 2,025 & 270 & 270 \\
              & UltraChat & 1,000 & 128 & 128 \\
              & \textbf{Total} & \textbf{9,025} & \textbf{782} & \textbf{782} \\
    \midrule
    DFlash & Nemotron V1 mathematics & 20,000 & 400 & -- \\
           & Nemotron V1 code / TACO train & 10,000 & 200 & -- \\
           & CodeAlpaca & 5,000 & 100 & -- \\
           & UltraChat & 15,000 & 300 & -- \\
           & \textbf{Total} & \textbf{50,000} & \textbf{1,000} & -- \\
    \bottomrule
  \end{tabular}
\end{table*}

\paragraph{Selection and teacher trajectories.}
We select prompts deterministically with seed 42 from fixed source
versions and shards.
We remove normalized exact duplicates within each training set and
screen candidates against evaluation questions, few-shot demonstrations,
and both methods' held-out splits using normalized matching, containment
checks, and token 5-gram similarity with a Jaccard threshold of 0.9.
These checks address lexical overlap but cannot exclude all semantic
paraphrases.
Each frozen target generates its own full-depth greedy responses,
which replace the source answers as training labels.
The maximum sequence length is 3,072 tokens, comprising at most
2,560 prompt tokens and 512 generated tokens,
using task-specific completion prompts for LoopAhead and the target's
released chat template for DFlash.
For DFlash, discarding responses shorter than 32 tokens leaves
48,137 training examples for Ouro-2.6B and 47,387 for Huginn-3.5B.
All 9,025 LoopAhead training prompts are retained for both targets.

\paragraph{LoopAhead configuration.}
The lookahead drafter has two causal Transformer layers with hidden
width 512, eight attention heads, and feed-forward width 1,536.
It conditions on 64 preceding target states and a teacher anchor token,
and learns to predict up to seven future tokens using cross-entropy.
States are extracted after the first recurrent pass for Ouro and after
the fourth pass, including the corresponding readout, for Huginn.
For Huginn, we train this base drafter directly on fourth-pass states
without the additional Raven-specific residual adaptation.
The resulting drafters contain 11.1M and 17.7M trainable parameters,
respectively.
We evaluate every 500 optimizer updates and at the final update,
selecting the checkpoint with the highest domain-averaged accepted
prefix length on the development set, with cross-entropy breaking ties.
The selected checkpoint is evaluated once on the confirmation set.

\paragraph{DFlash configuration.}
We use the five-layer block-diffusion drafter and training objective
implemented in SpecForge, with a block size of 16 comprising one anchor
and 15 masked positions. Five target features are fused and injected
into the drafter's key--value representations.
For Ouro, these features are taken from evenly spaced layers within
the final recurrent pass; for Huginn, they comprise the final prelude
layer, two layers of the final recurrent pass, and both coda layers.
For Ouro-2.6B, the drafter retains the target width and has 277.9M
trainable parameters. For Huginn-3.5B, we use a narrower drafter with
hidden width 2,048, feed-forward width 5,632, and 16 query and key--value
heads of dimension 128. Trainable linear projections map the frozen
5,280-dimensional target embeddings to the drafter width, fuse the five
target features from 26,400 to 2,048 dimensions, and map the drafter
output back to 5,280 dimensions for the frozen target output head.
This Huginn adaptation has 332,621,056 (332.6M) trainable parameters,
including the feature-fusion and input/output projections but excluding
the frozen target embeddings and output head, and is trained from scratch.
Training uses masked-token cross-entropy with exponential positional
decay of scale 7, and target features are computed online.
The evaluation protocol uses the final checkpoint after six epochs.

\paragraph{Optimization.}
Table~\ref{tab:drafter-optimization} lists the optimization settings.
Both methods use BF16 computation with FP32 optimizer master weights,
gradient-norm clipping at 1.0, and training seed 42 on NVIDIA L20X GPUs.
The two methods have distinct data and
optimization budgets, following their respective training recipes.

\begin{table}[t]
  \centering
  \small
  \setlength{\tabcolsep}{4pt}
  \caption{Optimization settings. Batch sizes are global.
  LoopAhead samples anchors per sequence per epoch; DFlash samples
  anchors each time a sequence is presented.}
  \label{tab:drafter-optimization}
  \begin{tabular}{@{}lcc@{}}
    \toprule
    Setting & LoopAhead & DFlash \\
    \midrule
    Epochs & 3 & 6 \\
    Batch size & 64 windows & 32 sequences \\
    Anchors per sequence & 32 & 512 \\
    Optimizer & AdamW & AdamW \\
    Adam $(\beta_1,\beta_2)$ & $(0.9,0.95)$ & $(0.9,0.999)$ \\
    Weight decay & 0.01 & 0 \\
    Peak learning rate & $2\times10^{-4}$ & $6\times10^{-4}$ \\
    Warmup fraction & 5\% & 4\% \\
    LR decay & Cosine & Cosine \\
    Minimum LR & $2\times10^{-5}$ & 0 \\
    \bottomrule
  \end{tabular}
\end{table}
